\section{Discussion}\label{sec:discussion}

\textbf{What the rare-event formulation provides.} The formulation in terms of $P(N)$ and $N_{0.9}$ turns the statement ``the method finds a hidden attractor'' into a measurable and comparable one. It makes visible three things that the qualitative setting obscures. First, many classical examples are not hard search problems, because the basin of the hidden attractor occupies a noticeable part of the natural region (Table~\ref{tab:benchmark}). Second, on truly rare targets the heuristics, including evolutionary search with Zelinka's cost function,\cite{Zelinka2016} are no better than uniform sampling and sometimes worse (Table~\ref{tab:gains}). Uniform sampling is itself a strong reference: its cost $2.3/p$ is known in advance, does not depend on the geometry of the basin, and is within a factor of 2.6 of the best possible label-only algorithm (Sec.~\ref{sec:channels}). Third, rarity is a property of the chosen region rather than of the attractor: enlarging the region or moving the parameter toward the crisis of the attractor lowers the basin measure by one or two orders of magnitude at most (Appendix~\ref{app:enlarged}). A truly rare hidden target in a real system arises from a combination of several hidden regimes, as in coupled units (Sec.~\ref{sec:coupled}).

\textbf{When adaptive search gains.} A gain is possible only through the information carried by trajectories that missed the target (Sec.~\ref{sec:channels}), and the data confirm this. Channel B (transients) gives an order-of-growth gain if the transient duration grows on approach to the basin on all scales, as in D2 and, judging by the diagnostic $d(q)$, in basins surrounded by edge states (Chua~II, LE1, SE11); the enlarged regions of Appendix~\ref{app:enlarged} give effective exponents of 0.4--0.7 on SE11, LE1, and Chua~I, with the caveat that the geometry of the basin changes with the region. The most favorable case is a system of weakly coupled units with hidden attractors: the transient duration reports which units have not yet settled, splitting fixes them one after another, and the gain grows geometrically with the number of units. The order-of-growth gain is bounded in $p$ by the time resolution of the coordinate, so the time step is a parameter of the search, not only of the integration. Channel A (boundaries) works if the boundary of the known basins is short and the target forms a noticeable part of it. Channel C (prior sets) gives a gain when the heuristic is fortunate (nullclines for Chua~I) and costs nothing when it is not. The diagnostic $d(q)$ separates the cases of gain and failure on all seven problems where it was measured, but it requires a known target and therefore serves as a tool of analysis rather than prediction. The splitting arm is a multichain relative of the stagger-and-step method,\cite{SweetNusseYorke2001} which climbs the escape time of single trajectories to locate chaotic saddles. A single-chain climb on the same coordinate is as efficient as splitting or more where the signal is present, including the coupled copies and D2 down to $p=10^{-5}$, and, without a uniform share, falls below the bound of the portfolio where the coordinate is deceptive or uninformative (Appendix~\ref{app:ablation}). What matters for the gain is the coordinate, not the particular way of climbing it, and a stagger-and-step chain is a natural arm of the portfolio, under the same guarantee.

\textbf{When the gain is a gain in time.} Counting integrations treats a trajectory that escapes after a few steps and one that runs over the whole horizon as equal. The method selects the long ones by construction, so the count overstates its advantage wherever most trajectories escape quickly, and there the gain in integrations is a loss or a draw in time (Sec.~\ref{sec:cost}). Where trajectories are bounded, as in dissipative physical systems, a nontarget integration is as expensive as a target one, and the gain carries over to steps and to time. In time, the gain also depends on how fast the reference evaluates a trajectory: against uniform sampling with recurrence-based stopping, which needs no observation window, it is 2--3 times smaller than against uniform sampling with our two-stage evaluation. The two choices are independent: the portfolio needs from each trajectory only a regime label and a transient duration, which the recurrence mapper provides, so the next step is to combine them. The method thus pays in time when nontarget integrations are not cheap compared with the long transients it selects, or when the target is so rare that the gain in integrations exceeds the ratio of their costs.

\textbf{Why a portfolio is needed.} In the worst cases, single adaptive strategies lose to uniform sampling several times and produce heavy tails of the hit-time distribution (Sec.~\ref{sec:failures}). By Proposition~\ref{prop:portfolio}, the portfolio with a uniform share $\theta$ is guaranteed to cost no more than $1/\theta$ times uniform sampling in integrations and inherits the order of the best autonomous arm. On all failure problems, the portfolio stayed within these limits, whereas splitting alone fell well below them (Sec.~\ref{sec:failures}, Appendix~\ref{app:enlarged}), and the uniform share costs little on the favorable problems (Appendix~\ref{app:ablation}). Bisection toward basin boundaries\cite{NusseYorke1989,Skufca2006Edge} is a strong baseline on systems with smooth boundaries and costs nothing where there is no signal, but it loses where the signal is graded over many scales and does not scale to coupled units (Appendix~\ref{app:ablation}); as a fourth arm, it helps only where it is strong on its own.

\textbf{Relation to continuation and to compositional search.} The established way to find a rare attractor is to continue it from parameter values where it is easy to find, by path following or by global continuation.\cite{Brzeski2017,Datseris2023Framework,Datseris2026Global} This requires such a path; our method works at fixed parameters, which is the situation when a new system is explored or when the attractor exists only in a narrow parameter window. The two approaches are complementary: attractors found at neighboring parameter values are natural prior sets for channel C, and the portfolio can replace the uniform sample within each continuation step. For coupled units whose individual attractors are known, product initial conditions are another prior set, and in the coupled Chua circuits they reach both the rare family and the coupling-born attractor directly; the portfolio is for the case when the units' attractors are not known or the regime is not a product.

\textbf{Verification as part of the method.} Three episodes of this work (Appendix~\ref{app:verification}) show that without a verification protocol a comparison of search methods loses its meaning. False ``hidden attractors'' due to transient chaos are known;\cite{SprottMunmuangsaen2018Comment,Danca2016Transient} false ``absences'' due to the integration step are, according to our data, just as real. The same protocol, applied to enumerate the regimes of a new system, found the coupling-born attractor of the Chua circuits (Sec.~\ref{sec:newattr}). The hiddenness measure $h_\varepsilon$ complements the binary definition; what matters in practice is the scale at which neighborhoods of the equilibria cease to lead to the attractor, and for rare targets the number of starts must be large enough for a zero to be informative (Sec.~\ref{sec:coupled}).

\textbf{Limitations.} All conclusions about gains refer to the benchmark, the synthetic families, and the coupled copies. On other systems, a gain is not guaranteed; only the worst-case bound of Proposition~\ref{prop:portfolio} is, and only in integrations. For maps and for systems with escape by slow drift, the reaction coordinate may be uninformative or deceptive, and in a large region, slowly returning trajectories make it deceptive even for Chua~II. The discretization limit bounds the dynamic range of the coordinate by the time step. The gains of $10^3$--$10^4$ in integrations are obtained on weakly coupled or uncoupled copies, whose rare target is close to a product of the units' attractors; in that setting, declaring the product states of the units known improves the gain, and the numbers without this prior knowledge are reported alongside (Sec.~\ref{sec:cost}). The same formulation applies to trained models of dynamics, such as reservoir computers in closed loop,\cite{Pathak2017Chaos,LuHuntOtt2018} which may develop attractors absent from the training data;\cite{HaluszczynskiRath2019,HerteuxRath2020,Flynn2021,Kabayama2025PseudoLorenz} that application will be reported separately.
